\chapter{Cauchy's Integral Formula}
\label{ch:iv06}

Cauchy's integral formula reconstructs a holomorphic function from boundary values. Its differentiated forms yield analyticity, derivative estimates, Liouville's theorem, and the mean-value property.

\section*{Learning goals}
After this chapter, the reader should be able to:
\begin{itemize}
\item apply Cauchy's integral formula to evaluate values and derivatives of holomorphic functions;
\item derive Cauchy estimates from contour bounds;
\item use Cauchy's formula to prove infinite differentiability and analyticity;
\item compute contour integrals containing powers of \((z-a)^{-1}\) by matching the correct derivative formula;
\item choose a circle contained in the domain to optimize an estimate;
\item use Cauchy estimates to prove Liouville-type or derivative-vanishing consequences;
\end{itemize}
\section*{Conceptual roadmap}
\[
\boxed{\text{local holomorphic structure}\;\longrightarrow\;\text{integral or series representation}\;\longrightarrow\;\text{global consequence}.}
\]

\section{Integral formula}
Inside a positively oriented simple closed contour, $f(z)$ equals a boundary integral with kernel $(\zeta-z)^{-1}$.

\section{Circle form}
On circles the formula becomes an explicit average weighted by the Cauchy kernel.

\section{Derivative formulas}
Differentiating under the integral gives formulas for every derivative.

\section{Cauchy estimates}
Derivative formulas bound derivatives by boundary maxima and powers of radius.

\section{Analyticity}
Holomorphic functions are automatically represented by convergent Taylor series.

\section{Mean-value property}
Holomorphic functions and their harmonic components satisfy circle-average identities.

\section{Liouville theorem}
A bounded entire function must be constant.

\section{Fundamental theorem of algebra}
Liouville's theorem yields a short proof that every nonconstant complex polynomial has a root.

\section{Core structural results}

\begin{theorem}[Cauchy integral formula]\label{thm:iv06-01}
If $f$ is holomorphic on and inside a positively oriented simple closed contour $\Gamma$ and $z$ lies inside, then $f(z)=\frac1{2\pi i}\int_\Gamma\frac{f(\zeta)}{\zeta-z}d\zeta$.
\begin{proof}
Subtract $f(z)$ from the numerator so the apparent singularity becomes removable, apply Cauchy's theorem to the regular part, and compute the integral of $(\zeta-z)^{-1}$ around the contour as $2\pi i$.
\end{proof}
\end{theorem}

\begin{theorem}[Cauchy derivative formula]\label{thm:iv06-02}
$f^{(n)}(z)=\frac{n!}{2\pi i}\int_\Gamma\frac{f(\zeta)}{(\zeta-z)^{n+1}}d\zeta$.
\begin{proof}
Differentiate the Cauchy formula with respect to $z$; uniform separation of $z$ from the contour justifies differentiation under the integral.
\end{proof}
\end{theorem}

\begin{theorem}[Cauchy estimate]\label{thm:iv06-03}
If $|f|\le M$ on $|\zeta-z_0|=R$, then $|f^{(n)}(z_0)|\le n!M/R^n$.
\begin{proof}
Apply the derivative formula on the circle and use the ML inequality with length $2\pi R$ and denominator magnitude $R^{n+1}$.
\end{proof}
\end{theorem}

\begin{theorem}[Liouville theorem]\label{thm:iv06-04}
Every bounded entire function is constant.
\begin{proof}
Apply the Cauchy estimate for the first derivative on circles of radius $R$ and let $R\to\infty$, forcing $f'=0$.
\end{proof}
\end{theorem}

\section{Worked examples}

\begin{example}[Integral formula diagnostic]\label{ex:iv06-worked-01}
Give a rigorous worked analysis of Integral formula. State the relevant holomorphic, contour, series, or singularity hypothesis and derive the central conclusion. Inside a positively oriented simple closed contour, $f(z)$ equals a boundary integral with kernel $(\zeta-z)^{-1}$. The decisive step is to use the complex-analytic structure globally rather than reason from real-variable intuition alone.
\end{example}

\begin{example}[Circle form diagnostic]\label{ex:iv06-worked-02}
Give a rigorous worked analysis of Circle form. State the relevant holomorphic, contour, series, or singularity hypothesis and derive the central conclusion. On circles the formula becomes an explicit average weighted by the Cauchy kernel. The decisive step is to use the complex-analytic structure globally rather than reason from real-variable intuition alone.
\end{example}

\begin{example}[Derivative formulas diagnostic]\label{ex:iv06-worked-03}
Give a rigorous worked analysis of Derivative formulas. State the relevant holomorphic, contour, series, or singularity hypothesis and derive the central conclusion. Differentiating under the integral gives formulas for every derivative. The decisive step is to use the complex-analytic structure globally rather than reason from real-variable intuition alone.
\end{example}

\begin{example}[Cauchy estimates diagnostic]\label{ex:iv06-worked-04}
Give a rigorous worked analysis of Cauchy estimates. State the relevant holomorphic, contour, series, or singularity hypothesis and derive the central conclusion. Derivative formulas bound derivatives by boundary maxima and powers of radius. The decisive step is to use the complex-analytic structure globally rather than reason from real-variable intuition alone.
\end{example}

\section{Solved dossiers}
The dossiers are canonical solved problems. The provenance ledger separately records which retained corpus rules guided each topic and which dossiers were newly designed.

\begin{problem}[Integral formula diagnostic]\label{prob:iv06-dossier-01}
Give a rigorous worked analysis of Integral formula. State the relevant holomorphic, contour, series, or singularity hypothesis and derive the central conclusion.
\end{problem}
\begin{solution}
Inside a positively oriented simple closed contour, $f(z)$ equals a boundary integral with kernel $(\zeta-z)^{-1}$. The decisive step is to use the complex-analytic structure globally rather than reason from real-variable intuition alone.
\end{solution}

\begin{problem}[Circle form diagnostic]\label{prob:iv06-dossier-02}
Give a rigorous worked analysis of Circle form. State the relevant holomorphic, contour, series, or singularity hypothesis and derive the central conclusion.
\end{problem}
\begin{solution}
On circles the formula becomes an explicit average weighted by the Cauchy kernel. The decisive step is to use the complex-analytic structure globally rather than reason from real-variable intuition alone.
\end{solution}

\begin{problem}[Derivative formulas diagnostic]\label{prob:iv06-dossier-03}
Give a rigorous worked analysis of Derivative formulas. State the relevant holomorphic, contour, series, or singularity hypothesis and derive the central conclusion.
\end{problem}
\begin{solution}
Differentiating under the integral gives formulas for every derivative. The decisive step is to use the complex-analytic structure globally rather than reason from real-variable intuition alone.
\end{solution}

\begin{problem}[Cauchy estimates diagnostic]\label{prob:iv06-dossier-04}
Give a rigorous worked analysis of Cauchy estimates. State the relevant holomorphic, contour, series, or singularity hypothesis and derive the central conclusion.
\end{problem}
\begin{solution}
Derivative formulas bound derivatives by boundary maxima and powers of radius. The decisive step is to use the complex-analytic structure globally rather than reason from real-variable intuition alone.
\end{solution}

\begin{problem}[Analyticity diagnostic]\label{prob:iv06-dossier-05}
Give a rigorous worked analysis of Analyticity. State the relevant holomorphic, contour, series, or singularity hypothesis and derive the central conclusion.
\end{problem}
\begin{solution}
Holomorphic functions are automatically represented by convergent Taylor series. The decisive step is to use the complex-analytic structure globally rather than reason from real-variable intuition alone.
\end{solution}

\begin{problem}[Mean-value property diagnostic]\label{prob:iv06-dossier-06}
Give a rigorous worked analysis of Mean-value property. State the relevant holomorphic, contour, series, or singularity hypothesis and derive the central conclusion.
\end{problem}
\begin{solution}
Holomorphic functions and their harmonic components satisfy circle-average identities. The decisive step is to use the complex-analytic structure globally rather than reason from real-variable intuition alone.
\end{solution}

\begin{problem}[Liouville theorem diagnostic]\label{prob:iv06-dossier-07}
Give a rigorous worked analysis of Liouville theorem. State the relevant holomorphic, contour, series, or singularity hypothesis and derive the central conclusion.
\end{problem}
\begin{solution}
A bounded entire function must be constant. The decisive step is to use the complex-analytic structure globally rather than reason from real-variable intuition alone.
\end{solution}

\begin{problem}[Cauchy integral formula proof dossier]\label{prob:iv06-dossier-08}
Prove the following theorem-level statement and identify the essential hypothesis: If $f$ is holomorphic on and inside a positively oriented simple closed contour $\Gamma$ and $z$ lies inside, then $f(z)=\frac1{2\pi i}\int_\Gamma\frac{f(\zeta)}{\zeta-z}d\zeta$.
\end{problem}
\begin{solution}
Subtract $f(z)$ from the numerator so the apparent singularity becomes removable, apply Cauchy's theorem to the regular part, and compute the integral of $(\zeta-z)^{-1}$ around the contour as $2\pi i$. This identifies the mechanism that makes the complex-analytic conclusion stronger than its real-variable analogue.
\end{solution}

\begin{problem}[Cauchy derivative formula proof dossier]\label{prob:iv06-dossier-09}
Prove the following theorem-level statement and identify the essential hypothesis: $f^{(n)}(z)=\frac{n!}{2\pi i}\int_\Gamma\frac{f(\zeta)}{(\zeta-z)^{n+1}}d\zeta$.
\end{problem}
\begin{solution}
Differentiate the Cauchy formula with respect to $z$; uniform separation of $z$ from the contour justifies differentiation under the integral. This identifies the mechanism that makes the complex-analytic conclusion stronger than its real-variable analogue.
\end{solution}

\begin{problem}[Cauchy estimate proof dossier]\label{prob:iv06-dossier-10}
Prove the following theorem-level statement and identify the essential hypothesis: If $|f|\le M$ on $|\zeta-z_0|=R$, then $|f^{(n)}(z_0)|\le n!M/R^n$.
\end{problem}
\begin{solution}
Apply the derivative formula on the circle and use the ML inequality with length $2\pi R$ and denominator magnitude $R^{n+1}$. This identifies the mechanism that makes the complex-analytic conclusion stronger than its real-variable analogue.
\end{solution}

\begin{problem}[Liouville theorem proof dossier]\label{prob:iv06-dossier-11}
Prove the following theorem-level statement and identify the essential hypothesis: Every bounded entire function is constant.
\end{problem}
\begin{solution}
Apply the Cauchy estimate for the first derivative on circles of radius $R$ and let $R\to\infty$, forcing $f'=0$. This identifies the mechanism that makes the complex-analytic conclusion stronger than its real-variable analogue.
\end{solution}

\begin{problem}[Chapter synthesis]\label{prob:iv06-dossier-12}
Connect the definitions and principal theorems of this chapter into one reusable complex-analysis strategy.
\end{problem}
\begin{solution}
Cauchy's integral formula reconstructs a holomorphic function from boundary values. Its differentiated forms yield analyticity, derivative estimates, Liouville's theorem, and the mean-value property. The reusable strategy is to identify the analytic domain, choose the correct local expansion or contour representation, apply the structural theorem, and then audit singularities and topology.
\end{solution}

\section{Exercises with complete solutions}

\begin{exercise}\label{exr:iv06-01}
State and apply the central principle behind Integral formula. Give one concrete consequence.
\end{exercise}
\begin{hint}
Match the integrand to \(f(z)/(z-a)\) and check that \(a\) lies inside the positively oriented contour.
\end{hint}
\begin{solution}
Inside a positively oriented simple closed contour, $f(z)$ equals a boundary integral with kernel $(\zeta-z)^{-1}$. The same principle can then be reused in later contour, residue, conformal, or Riemann-surface arguments.
\end{solution}

\begin{exercise}\label{exr:iv06-02}
State and apply the central principle behind Circle form. Give one concrete consequence.
\end{exercise}
\begin{hint}
For higher derivatives, rewrite the denominator as \((z-a)^{n+1}\) and use \(2\pi i f^{(n)}(a)/n!\).
\end{hint}
\begin{solution}
On circles the formula becomes an explicit average weighted by the Cauchy kernel. The same principle can then be reused in later contour, residue, conformal, or Riemann-surface arguments.
\end{solution}

\begin{exercise}\label{exr:iv06-03}
State and apply the central principle behind Derivative formulas. Give one concrete consequence.
\end{exercise}
\begin{hint}
Choose a circle centered at \(a\) contained in the domain, then bound the Cauchy integral by its maximum modulus.
\end{hint}
\begin{solution}
Differentiating under the integral gives formulas for every derivative. The same principle can then be reused in later contour, residue, conformal, or Riemann-surface arguments.
\end{solution}

\begin{exercise}\label{exr:iv06-04}
State and apply the central principle behind Cauchy estimates. Give one concrete consequence.
\end{exercise}
\begin{hint}
Cauchy's estimate gives \(|f^{(n)}(a)|\le n!M/R^n\); select \(R\) to make the bound useful.
\end{hint}
\begin{solution}
Derivative formulas bound derivatives by boundary maxima and powers of radius. The same principle can then be reused in later contour, residue, conformal, or Riemann-surface arguments.
\end{solution}

\begin{exercise}\label{exr:iv06-05}
State and apply the central principle behind Analyticity. Give one concrete consequence.
\end{exercise}
\begin{hint}
To prove analyticity, expand \(1/(z-w)\) geometrically for \(|w-a|<|z-a|\) inside the Cauchy formula.
\end{hint}
\begin{solution}
Holomorphic functions are automatically represented by convergent Taylor series. The same principle can then be reused in later contour, residue, conformal, or Riemann-surface arguments.
\end{solution}

\begin{exercise}\label{exr:iv06-06}
State and apply the central principle behind Mean-value property. Give one concrete consequence.
\end{exercise}
\begin{hint}
For Liouville-type reasoning, let the radius in a Cauchy estimate tend to infinity while the global bound remains fixed.
\end{hint}
\begin{solution}
Holomorphic functions and their harmonic components satisfy circle-average identities. The same principle can then be reused in later contour, residue, conformal, or Riemann-surface arguments.
\end{solution}

\begin{exercise}\label{exr:iv06-07}
State and apply the central principle behind Liouville theorem. Give one concrete consequence.
\end{exercise}
\begin{hint}
If the contour does not wind once around \(a\), include the winding number rather than using the simple formula blindly.
\end{hint}
\begin{solution}
A bounded entire function must be constant. The same principle can then be reused in later contour, residue, conformal, or Riemann-surface arguments.
\end{solution}

\begin{exercise}\label{exr:iv06-08}
State and apply the central principle behind Fundamental theorem of algebra. Give one concrete consequence.
\end{exercise}
\begin{hint}
Audit whether singularities of the auxiliary function lie inside the chosen contour before differentiating under the formula.
\end{hint}
\begin{solution}
Liouville's theorem yields a short proof that every nonconstant complex polynomial has a root. The same principle can then be reused in later contour, residue, conformal, or Riemann-surface arguments.
\end{solution}

\section*{Chapter summary}
The chapter now forms part of the canonical complex-analysis chain from holomorphicity through residues, global continuation, and Riemann surfaces.
